% Arithmetic-geometry layer for the fold gauge-anomaly spectral curve.

\paragraph{谱曲线的算术几何与分歧多项式的同域结构}\label{par:fold-gauge-anomaly-spectral-curve-arithmetic}
命题 \ref{prop:fold-gauge-anomaly-discriminant-factorization} 给出谱四次 \(F(\mu,u)\) 的判别式分解与分支集合，
命题 \ref{prop:fold-gauge-anomaly-rate-curve-param} 给出速率曲线的消元模型 \(R(\alpha,u)=0\) 及其可行分支选择。
基于同一组代数对象，可以进一步抽取出谱曲线、分歧多项式与速率曲线之间的几何与算术对齐关系。

\begin{proposition}[谱四次曲线的平面模型、亏格与简单分歧]\label{prop:fold-gauge-anomaly-spectral-quartic-geometry}
沿用命题 \ref{prop:fold-gauge-anomaly-discriminant-factorization} 的记号。
令 \(F^{\mathrm{hom}}(\mu,u,w)\) 为 \(F(\mu,u)\) 的总次数齐次化，并记平面曲线
\[
C:=\left\{[\mu:u:w]\in \PP^{2}:\ F^{\mathrm{hom}}(\mu,u,w)=0\right\},
\]
其中
\[
F^{\mathrm{hom}}(\mu,u,w)
=\mu^{4}-\mu^{3}w-2\mu^{2}uw-\mu^{2}w^{2}-\mu u^{3}+2\mu u w^{2}+2u w^{3}.
\]
则 \(C\) 为光滑平面四次曲线，因而其亏格为 \(3\)，并且为非超椭圆曲线。
投影
\[
\pi:C\to\PP^{1},\qquad [\mu:u:w]\mapsto [u:w],
\]
为次数 \(4\) 的覆盖；其分支集合由 \(\mathrm{Disc}_{\mu}(F(\mu,u))=0\) 给出，等价地由
\(
u(u-1)P_{10}(u)=0
\)
支撑。
并且在全部分支点处均为单分歧（分歧指数为 \(2\)）。
\end{proposition}

\begin{proposition}[分歧多项式 \texorpdfstring{$P_{10}$}{P10} 的满对称伽罗瓦群与判别式]\label{prop:fold-gauge-anomaly-P10-galois-discriminant}
令 \(P_{10}(u)\in\ZZ[u]\) 如命题 \ref{prop:fold-gauge-anomaly-discriminant-factorization}。
则其伽罗瓦群为满对称群
\[
\mathrm{Gal}\bigl(P_{10}/\QQ\bigr)\cong S_{10},
\]
并且判别式具有显式分解
\[
\boxed{\ \mathrm{Disc}\bigl(P_{10}\bigr)=2^{38}3^{24}5^{2}\cdot 1931\cdot 9013^{2}\cdot 34319^{3}\ }.
\]
因此，除去平凡分支 \(u\in\{0,1\}\) 后，谱四次曲线的非平凡分歧位置在 \(\QQ\) 上达到满对称的算术一般性。
\end{proposition}

\begin{proposition}[重根谱值多项式与 \texorpdfstring{$P_{10}$}{P10} 的 Tschirnhaus 同域同构]\label{prop:fold-gauge-anomaly-Q10-tschirnhaus}
考虑系统
\[
F(\mu,u)=0,\qquad \partial_{\mu}F(\mu,u)=0.
\]
消去 \(u\) 并剔除 \(u\in\{0,1\}\) 的平凡分歧贡献后，得到关于 \(\mu\) 的次数 \(10\) 多项式
\[
\begin{aligned}
Q_{10}(\mu)=\;&27\mu^{10}-81\mu^{9}+90\mu^{8}-46\mu^{7}+11\mu^{6}-\mu^{5}\\
&-32\mu^{4}+32\mu^{3}+16\mu+16,
\end{aligned}
\]
其根精确刻画所有非平凡分支点处的重根谱值（即同时满足 \(F=\partial_{\mu}F=0\) 的 \(\mu\)-坐标）。
并且
\[
\mathrm{Gal}\bigl(Q_{10}/\QQ\bigr)\cong S_{10},\qquad
\mathrm{Disc}\bigl(Q_{10}\bigr)=2^{42}3^{16}5^{8}\cdot 1931\cdot 34319.
\]
更强地，\(P_{10}\) 与 \(Q_{10}\) 生成同一个次数 \(10\) 数域：存在显式的 Tschirnhaus 变换
\[
\begin{aligned}
u=\frac{1}{200}\bigl(&-189\mu^{9}+540\mu^{8}-360\mu^{7}-308\mu^{6}+329\mu^{5}\\
&+304\mu^{4}-104\mu^{3}-196\mu^{2}-128\mu-16\bigr),
\end{aligned}
\]
使得对任意满足 \(Q_{10}(\mu)=0\) 的根，上式给出的 \(u\) 自动满足 \(P_{10}(u)=0\)。
因此分歧位置的数域与谱碰撞值的数域在算术意义上完全识别。
\end{proposition}

\begin{corollary}[判别式平方比值的同域证书]\label{cor:fold-gauge-anomaly-disc-square-ratio}
在命题 \ref{prop:fold-gauge-anomaly-P10-galois-discriminant} 与命题 \ref{prop:fold-gauge-anomaly-Q10-tschirnhaus} 的记号下，
\[
\boxed{\ \frac{\mathrm{Disc}(P_{10})}{\mathrm{Disc}(Q_{10})}=\Bigl(\frac{25054688907}{500}\Bigr)^{2}\ }.
\]
因此，即使不显式写出极大整环的整基，上式仍在纯判别式层面给出同域证书：两者所诱导的整环阶之差异仅体现在指数的平方因子上。
\end{corollary}

\begin{proposition}[速率曲线消元模型的判别式平方因子分解]\label{prop:fold-gauge-anomaly-rate-curve-discriminant-squarefactor}
令 \(R(\alpha,u)\in\ZZ[\alpha,u]\) 为命题 \ref{prop:fold-gauge-anomaly-rate-curve-param} 的消元证书（表 \ref{tab:fold_gauge_anomaly_rate_curve_resultant_degree}），其定义使得速率曲线由 \(R(\alpha,u)=0\) 刻画且 \(\deg_{\alpha}R=4\)。
则其关于 \(\alpha\) 的判别式满足精确分解
\[
\boxed{\ \mathrm{Disc}_{\alpha}\bigl(R(\alpha,u)\bigr)=-(u^{5}H(u))^{2}\cdot u(u-1)P_{10}(u)\ },
\]
其中 \(H(u)\in\ZZ[u]\) 为次数 \(19\) 的不可约多项式
\[
\begin{aligned}
H(u)=\;&137781u^{19}+629856u^{18}+934578u^{17}-1546209u^{16}-6187752u^{15}+1402704u^{14}\\
&+15349014u^{13}-3368736u^{12}-17688806u^{11}+2216732u^{10}+17425320u^{9}-4765670u^{8}\\
&-11620472u^{7}+7448180u^{6}+2529720u^{5}-4736240u^{4}+2386300u^{3}-612800u^{2}+81800u-4500.
\end{aligned}
\]
因此，速率曲线作为 \(u\)-直线上的次数 \(4\) 覆盖，其平方自由判别式与谱四次曲线的平方自由判别式一致，分歧支撑同为 \(u(u-1)P_{10}(u)=0\)；
而 \(H(u)^{2}\) 以纯平方因子出现，反映了消元坐标环与其整闭包之间的非平凡指数现象，而非分歧的消失。
\end{proposition}

\begin{proposition}[指数平方因子的算术一般性：\texorpdfstring{$H$}{H} 的满对称伽罗瓦群与判别式]\label{prop:fold-gauge-anomaly-H-galois-discriminant}
令 \(H(u)\) 如命题 \ref{prop:fold-gauge-anomaly-rate-curve-discriminant-squarefactor}。
则
\[
\mathrm{Gal}\bigl(H/\QQ\bigr)\cong S_{19},
\]
并且其判别式分解为
\[
\begin{aligned}
\mathrm{Disc}(H)=\;&2^{40}3^{91}5^{32}\cdot 101\cdot 173^{2}\cdot 1931^{10}\cdot 9013^{3}\\
&\cdot 795449^{6}\cdot 15208978734786979108724207.
\end{aligned}
\]
该性质表明：\(\mathrm{Disc}_{\alpha}(R)\) 中的平方因子并非来自低阶可解对称性，而在伽罗瓦意义上同样处于满对称的一般情形。
\end{proposition}

\begin{proposition}[谱四次的三次预解式与判别式恒等]\label{prop:fold-gauge-anomaly-cubic-resolvent-discriminant}
对谱四次 \(F(\mu,u)\) 取标准三次预解式，得到
\[
\mathcal{R}(y,u)=y^{3}+(1+2u)y^{2}+(u^{3}-10u)y-\bigl(u^{6}-4u^{4}+20u^{2}+10u\bigr).
\]
则存在判别式恒等式
\[
\boxed{\ \mathrm{Disc}_{y}\bigl(\mathcal{R}(y,u)\bigr)=\mathrm{Disc}_{\mu}\bigl(F(\mu,u)\bigr)=-u(u-1)P_{10}(u)\ }.
\]
因此谱四次的 \(S_{4}\) 伽罗瓦闭包在 \(\QQ(u)\) 上携带由 \(\mathcal{R}\) 显式给出的规范 \(S_{3}\) 中间扩张，其分歧位置与谱四次曲线完全一致。
\end{proposition}

\begin{proposition}[一参数变形的分歧共振唯一性与泛性伽罗瓦群]\label{prop:fold-gauge-anomaly-branch-resonance-one-parameter}
将命题 \ref{prop:fold-gauge-anomaly-pressure} 中矩阵 \(2A_{\theta}\) 的 \((3,2)\) 元从 \(2\) 变形为参数 \(t\)，得到谱多项式
\[
F_{t}(\mu,u)=\mu^{4}-\mu^{3}-(tu+1)\mu^{2}+\mu(tu-u^{3})+tu.
\]
其关于 \(\mu\) 的判别式可写为
\[
\mathrm{Disc}_{\mu}\bigl(F_{t}(\mu,u)\bigr)=u\cdot D_{t}(u),
\]
其中 \(D_{t}(u)\in\ZZ[t][u]\) 为关于 \(u\) 的次数 \(11\) 多项式。
并且满足以下刚性算术性质：
\begin{enumerate}
\item 在实参数 \(t\in\RR\) 中，\(u=1\) 为分支点（等价于 \(D_{t}(1)=0\)）的唯一取值为 \(t=2\)；且在该点出现分解
\[
\boxed{\ D_{2}(u)=-(u-1)P_{10}(u)\ }.
\]
\item 在 \(t=1\) 时，\(D_{1}(u)\) 在 \(\QQ\) 上不可约且
\(
\mathrm{Gal}\bigl(D_{1}/\QQ\bigr)\cong S_{11}
\)；
由 Hilbert 不可约性可推出该一参数族在 \(\QQ(t)\) 上的通用分歧多项式具有泛性伽罗瓦群 \(S_{11}\)。
\end{enumerate}
因此 \(t=2\) 情形出现的异常可约分歧具有严格的唯一性框架：它是该一参数族中的唯一实分歧共振点。
\end{proposition}

\paragraph{单值群、符号二次商与 \texorpdfstring{$S_4$}{S4}-Hurwitz 塔}\label{par:fold-gauge-anomaly-s4-hurwitz-tower}
以下把命题 \ref{prop:fold-gauge-anomaly-spectral-quartic-geometry}--\ref{prop:fold-gauge-anomaly-cubic-resolvent-discriminant} 中给出的“谱四次—判别式—三次预解式”数据，转写为覆叠的 Hurwitz 型、泛型伽罗瓦群以及关键中间商曲线的亏格。
为与外部记号对齐，记
\[
D_{10}(u):=P_{10}(u),\qquad D_{19}(u):=H(u).
\]

\begin{conclusion}[规范异常四次谱曲线的几何型]\label{con:fold-gauge-anomaly-spectral-quartic-geomtype}
令 \(\mathcal C_F\subset\PP^2\) 为谱四次 \(F(\mu,u)=0\) 的总次数齐次化闭包（命题 \ref{prop:fold-gauge-anomaly-spectral-quartic-geometry} 中的曲线 \(C\)）。
则 \(\mathcal C_F\) 在特征 \(0\) 上无奇点，因而为非超椭圆的平面四次曲线，并满足
\[
g(\mathcal C_F)=3.
\]
\end{conclusion}

\begin{conclusion}[判别式的完全分解与简单分支]\label{con:fold-gauge-anomaly-disc-factorization-simple-branch}
将 \(F(\mu,u)\) 视为关于 \(\mu\) 的四次多项式，则其判别式满足
\[
\mathrm{Disc}_{\mu}\bigl(F(\mu,u)\bigr)=-u(u-1)D_{10}(u),
\qquad D_{10}(u)=27u^{10}+27u^{9}-153u^{8}-163u^{7}+793u^{6}+1061u^{5}-832u^{4}-816u^{3}+1320u^{2}-440u+40,
\]
并且 \(u(u-1)D_{10}(u)\) 在 \(\ZZ[u]\) 中平方因子自由。
因此投影 \(\pi:\mathcal C_F\to\PP^1_u\) 为次数 \(4\) 覆叠，其分支集合恰为
\[
\{u=0\}\cup\{u=1\}\cup\{u:\ D_{10}(u)=0\},
\]
且在每个分支值之上均为简单分支（惯性型为单一换位）。
\end{conclusion}
\begin{proof}[证明要点]
判别式因式分解见命题 \ref{prop:fold-gauge-anomaly-discriminant-factorization}。
平方因子自由性等价于
\(
\gcd\!\bigl(D_{10},D_{10}'\bigr)=\gcd(D_{10},u)=\gcd(D_{10},u-1)=1
\)
在 \(\ZZ[u]\) 中成立；从而各分支值均对应 \(\mu\)-纤维中的二重并合且无更高阶并合，故为简单分支。
\end{proof}

\begin{conclusion}[特征多项式的泛型伽罗瓦群为 \texorpdfstring{$S_4$}{S4}]\label{con:fold-gauge-anomaly-generic-galois-s4}
函数域扩张 \(\QQ(u)\subset \QQ(u,\mu)\)（由 \(F(\mu,u)=0\) 给出）之泛型伽罗瓦群满足
\[
\mathrm{Gal}\bigl(F(\mu,u)/\QQ(u)\bigr)\cong S_4.
\]
等价地，覆叠 \(\pi:\mathcal C_F\to\PP^1_u\) 的几何单值群为满对称群 \(S_4\)。
\end{conclusion}
\begin{proof}[证明要点（可核验特化证书）]
由命题 \ref{prop:fold-gauge-anomaly-cubic-resolvent-discriminant}，谱四次的三次预解式 \(\mathcal R(y,u)\) 满足
\(
\mathrm{Disc}_y(\mathcal R)=\mathrm{Disc}_{\mu}(F)=-u(u-1)D_{10}(u)
\)，其在 \(\QQ(u)\) 中非平方，故伽罗瓦群不包含于 \(A_4\)。
取 \(u=2\) 的特化，四次多项式 \(F(\mu,2)\in\ZZ[\mu]\) 在模 \(3\) 下不可约，而在模 \(11\) 下分解型为 \((3,1)\)，从而特化伽罗瓦群含有 \(3\)-循环并因此包含 \(A_4\)。
合并得特化伽罗瓦群为 \(S_4\)，从而泛型伽罗瓦群亦为 \(S_4\)。
\end{proof}

\begin{conclusion}[符号二次子覆盖：判别式双覆盖的亏格]\label{con:fold-gauge-anomaly-sign-quadratic-cover}
对应于 \(\mathrm{Gal}(F/\QQ(u))\cong S_4\) 的符号表示之二次子域由 \(\sqrt{\mathrm{Disc}_{\mu}(F)}\) 生成，其几何模型可取为双覆盖
\[
\mathcal H:\quad w^2=-u(u-1)D_{10}(u).
\]
由于右端次数为 \(12\)，故 \(g(\mathcal H)=5\)，并且 \(\mathcal H\to\PP^1_u\) 的分支集合与结论 \ref{con:fold-gauge-anomaly-disc-factorization-simple-branch} 完全一致。
\end{conclusion}

\begin{conclusion}[显式三次分辨：同一判别式核]\label{con:fold-gauge-anomaly-cubic-resolvent}
与 \(F(\mu,u)\) 相关联的三次预解式（命题 \ref{prop:fold-gauge-anomaly-cubic-resolvent-discriminant} 中的 \(\mathcal R(y,u)\)）可记为
\[
\Phi(y,u):=\mathcal R(y,u)=y^3+(1+2u)y^2+(u^3-10u)y-\bigl(u^6-4u^4+20u^2+10u\bigr),
\]
并满足判别式恒等式
\[
\mathrm{Disc}_{y}\bigl(\Phi(y,u)\bigr)=-u(u-1)D_{10}(u).
\]
因此 \(\Phi(y,u)=0\) 在 \(\PP^1_u\) 上给出一个次数 \(3\) 的不可约覆叠，其分支参数与 \(\pi:\mathcal C_F\to\PP^1_u\) 完全同位，且在每一分支值处呈现简单分支。
\end{conclusion}

\begin{conclusion}[本征谱的实–复相图由 \texorpdfstring{$D_{10}$}{D10} 的符号刻画]\label{con:fold-gauge-anomaly-real-complex-phase}
设 \(u>0\) 为实参数，并将 \(F(\mu,u)\) 在 \(\RR[\mu]\) 中的根型按实根/复根计数分类。
则 \(D_{10}(u)=0\) 在区间 \((0,1)\) 内恰有两条实根
\(
0<u_-<u_+<1
\)，并且：
\begin{enumerate}
  \item 当 \(0<u<u_-\) 或 \(u_+<u<1\) 时，\(F(\mu,u)\) 具有四条互异实根；
  \item 当 \(u_-<u<u_+\) 时，\(F(\mu,u)\) 具有两条互异实根与一对共轭非实复根；
  \item 当 \(u>1\) 时，\(F(\mu,u)\) 恒具有两条互异实根与一对共轭非实复根。
\end{enumerate}
此外，\(D_{10}(u)\) 在 \((-\infty,0)\cup(1,\infty)\) 无实根，故上述相图不存在其他实临界点。
\end{conclusion}
\begin{proof}[证明要点]
由结论 \ref{con:fold-gauge-anomaly-disc-factorization-simple-branch} 与平方因子自由性可知：除去 \(\{0,1\}\cup\{D_{10}=0\}\) 外，根型在各连通区间内常值。
又由于在 \(u\in(0,1)\) 上有 \(\mathrm{Disc}_{\mu}(F)=u(1-u)D_{10}(u)\)，故判别式符号等价于 \(D_{10}\) 的符号。
选取代表点（例如 \(u=\frac1{20},\frac15,\frac12,\frac65\)）作一次实根计数核验即可确定三类区间的根型；而 \(D_{10}\) 在 \((-\infty,0)\cup(1,\infty)\) 无实根见命题 \ref{prop:fold-gauge-anomaly-branch-point-classification} 的根枚举。
\end{proof}

\begin{conclusion}[速率曲线判别式分解与新的 \texorpdfstring{$19$}{19} 次平方因子]\label{con:fold-gauge-anomaly-rate-curve-discriminant-D19}
将消元证书 \(R(\alpha,u)\) 视为关于 \(\alpha\) 的四次多项式，则其判别式满足
\[
\mathrm{Disc}_{\alpha}\bigl(R(\alpha,u)\bigr)= -u^{11}(u-1)D_{10}(u)\,D_{19}(u)^2,
\]
其中 \(D_{19}(u)\in\ZZ[u]\) 可取为命题 \ref{prop:fold-gauge-anomaly-rate-curve-discriminant-squarefactor} 中的 \(H(u)\)。
并且 \(D_{19}(u)\) 在实轴上恰有 \(7\) 条实根，且全部落于区间 \((0,1)\)；在 \((-\infty,0)\cup(1,\infty)\) 无实根。
因此 \(D_{19}(u)^2\) 仅以纯平方因子形式进入判别式分解：它刻画了消元坐标环的指数型退化位形，而不改变速率曲线覆叠在平方自由意义下的分支支撑 \(\{u=0,1\}\cup\{D_{10}(u)=0\}\)。
\end{conclusion}
\begin{proof}[证明要点]
判别式分解见命题 \ref{prop:fold-gauge-anomaly-rate-curve-discriminant-squarefactor} 并注意 \((u^5D_{19}(u))^2=u^{10}D_{19}(u)^2\)。
实根计数与区间定位可由 Sturm 计数或高精度根枚举复核。
\end{proof}

\begin{conclusion}[速率曲线的泛型单值群同为 \texorpdfstring{$S_4$}{S4}]\label{con:fold-gauge-anomaly-rate-curve-generic-galois-s4}
域扩张 \(\QQ(u)\subset\QQ(u,\alpha)\)（由 \(R(\alpha,u)=0\) 定义）之泛型伽罗瓦群满足
\[
\mathrm{Gal}\bigl(R(\alpha,u)/\QQ(u)\bigr)\cong S_4.
\]
\end{conclusion}
\begin{proof}[证明要点（可核验特化证书）]
取 \(u=2\) 的特化并由结果式消元得到一元四次多项式 \(R(\alpha,2)\in\ZZ[\alpha]\)。
该多项式在模 \(3\) 下不可约且在模 \(11\) 下分解型为 \((3,1)\)，并且判别式非平方；
故其特化伽罗瓦群为 \(S_4\)，从而泛型亦为 \(S_4\)。
\end{proof}

\begin{conclusion}[\texorpdfstring{$\mathcal C_F$}{CF} 的无穷远有理点与三重覆盖结构]\label{con:fold-gauge-anomaly-trigonal-structure}
曲线 \(\mathcal C_F\subset\PP^2_{[\mu:u:w]}\) 在无穷远直线 \(w=0\) 上恰含两点 \(\QQ\)-有理点
\[
P_{\infty}^{(0)}=[0:1:0],\qquad P_{\infty}^{(1)}=[1:1:0],
\]
且二者皆为光滑点。
因此，从任一 \(P_{\infty}^{(i)}\) 作线性投影均可在 \(\QQ\) 上诱导一个次数 \(3\) 的态射
\(
\varphi_i:\mathcal C_F\to\PP^1
\)，从而 \(\mathcal C_F\) 具有 \(\QQ\)-三重覆盖结构（trigonal structure）。
\end{conclusion}
\begin{proof}[证明要点]
由齐次方程 \(F^{\mathrm{hom}}(\mu,u,w)=0\) 在 \(w=0\) 上化为 \(\mu(\mu^3-u^3)=0\) 可得无穷远点集合；
其中 \(\mu=0\) 与 \(\mu=u\) 给出上述两点的有理性。
又由 \(\partial_{\mu}F^{\mathrm{hom}}|_{w=0}=4\mu^3-u^3\) 在两点处非零知其为光滑点。
线性投影给出的线系在光滑点处可延拓为正则态射，且对平面四次曲线其次数为 \(3\)。
\end{proof}

\begin{conclusion}[三次分辨曲线的平面六次模型与正规化亏格]\label{con:fold-gauge-anomaly-resolvent-sextic-normalization-genus}
令 \(\mathcal C_3\subset\PP^2_{[y:u:w]}\) 为 \(\Phi(y,u)=0\) 的射影闭包。
则 \(\mathcal C_3\) 为一条平面六次曲线，且无穷远处仅含一点 \([1:0:0]\)；其仿射部分为光滑曲线，而 \([1:0:0]\) 为唯一奇点。
记 \(\widetilde{\mathcal C}_3\) 为其正规化，则
\[
g(\widetilde{\mathcal C}_3)=4,
\]
并且覆叠 \(\widetilde{\mathcal C}_3\to\PP^1_u\) 为次数 \(3\) 的映射，在结论 \ref{con:fold-gauge-anomaly-disc-factorization-simple-branch} 的 \(12\) 个分支参数处均为简单分支，而在 \(u=\infty\) 处不分支。
等价地，\([1:0:0]\) 的 \(\delta\)-不变量为 \(6\)。
\end{conclusion}
\begin{proof}[证明要点]
齐次化可见 \(\mathcal C_3\) 的无穷远点满足 \(u=0=w\)，故唯一为 \([1:0:0]\)。
由于 \(\mathrm{Disc}_y(\Phi)=-u(u-1)D_{10}(u)\) 且平方因子自由，有限分支值均为简单分支，故除无穷远处外曲线无奇点。
进一步，由 \(u\to\infty\) 时的主导平衡 \(y^3-u^6=0\) 得到三条互异分支 \(y\sim \zeta u^2\)（\(\zeta^3=1\)），从而 \(u=\infty\) 上方无分歧。
因此 \(\widetilde{\mathcal C}_3\to\PP^1_u\) 为次数 \(3\) 的覆叠且仅在 \(12\) 个有限点简单分支，
由 Riemann--Hurwitz 公式得 \(g(\widetilde{\mathcal C}_3)=4\)。
平面六次的算术亏格为 \(10\)，故 \(\delta=10-4=6\)。
\end{proof}

\begin{conclusion}[\texorpdfstring{$S_4$}{S4}-Hurwitz 塔：伽罗瓦闭包与中间商曲线的亏格]\label{con:fold-gauge-anomaly-s4-hurwitz-tower-genus}
令 \(\mathcal X\to\PP^1_u\) 表示由谱四次 \(F(\mu,u)=0\) 的 \(S_4\) 伽罗瓦闭包所对应的正规伽罗瓦覆叠（次数 \(24\)，群 \(S_4\)），并以结论 \ref{con:fold-gauge-anomaly-disc-factorization-simple-branch} 的 \(12\) 个简单分支点作为全部分支数据。
则 \(\mathcal X\) 的亏格为
\[
g(\mathcal X)=49.
\]
并且对以下中间子群的商曲线具有完全显式的亏格分层：
\begin{enumerate}
  \item 对 \(A_4\) 的商曲线同构于结论 \ref{con:fold-gauge-anomaly-sign-quadratic-cover} 的判别式双覆盖 \(\mathcal H\)，故亏格为 \(5\)；
  \item 对任一根稳定子 \(S_3\) 的商曲线同构于 \(\mathcal C_F\)，故亏格为 \(3\)；
  \item 对任一稳定配对结构的二面体子群 \(D_8\)（阶 \(8\)）的商曲线同构于结论 \ref{con:fold-gauge-anomaly-resolvent-sextic-normalization-genus} 的正规化 \(\widetilde{\mathcal C}_3\)，故亏格为 \(4\)；
  \item 对正规 Klein 四子群 \(V_4\) 的商覆叠次数为 \(6\)，且每一分支值的惯性在该 \(6\) 叶作用下为 \((2,2,2)\)，从而该商曲线的亏格为 \(13\)。
\end{enumerate}
\end{conclusion}
\begin{proof}[证明要点]
在 \(S_4\) 伽罗瓦覆叠中，每个分支点的惯性群由换位生成，故惯性阶为 \(2\)。
由 Riemann--Hurwitz 公式，
\(
2g(\mathcal X)-2=|S_4|\cdot(-2)+12\cdot|S_4|\,(1-\tfrac12)
\)
立得 \(g(\mathcal X)=49\)。
各中间商的亏格由同一公式对相应的置换表示（等价于对商覆叠的次数与惯性分解型）逐一计算得到；
其中 \(S_4/D_8\) 给出三次预解式的根域覆盖，而 \(S_4/V_4\cong S_3\) 给出其 \(S_3\) 伽罗瓦闭包。
\end{proof}

\begin{remark}[计算审计来源]\label{rem:fold-gauge-anomaly-spectral-curve-arithmetic-audit}
上述伽罗瓦群判定、判别式因子分解与 Tschirnhaus 变换系数可由计算代数数论的标准算法复现；其审计记录见 \cite{Internal2026FoldGaugeAnomalySpectralCurveArithmeticAudit20260218}。
\end{remark}

